\medskip

\begin{enumerate}
\item  On a : $z_1 = \sqrt{3} - \i$.

Son module vaut : $|z_1|=\sqrt{(\sqrt{3})^2+(-1)^2}=\sqrt{3+1}=2$.

On cherche un argument $\theta$ de $z_1$.

On a $\cos(\theta)=\dfrac{\sqrt{3}}{2}$ et $\sin(\theta)=-\dfrac{1}{2}$, on peut donc prendre $\theta=-\dfrac{\pi}{6}$.

Ainsi, la forme exponentielle de $z_1$ est : $z_1=2\e^{-\i\frac{\pi}{6}}$.

\item D'après la question précédente, $z_1=2\e^{-\i\frac{\pi}{6}}$ et, d'après l'énoncé, $z_2=2\e^{\i\frac{\pi}{4}}$. Donc:
\[
\begin{aligned}
z&=z_1^3\times z_2^2\\
&=\left(2\e^{-\i\frac{\pi}{6}}\right)^3\times\left(2\e^{\i\frac{\pi}{4}}\right)^2\\
&=8\e^{-\i\frac{\pi}{2}}\times 4\e^{\i\frac{\pi}{2}}\\
&=32\e^0\\
&=32.
\end{aligned}
\]
Le nombre $z$ est donc bien un réel.
\end{enumerate}

\bigskip
